% ==============================================================================
% SECCION 7: CUESTIONARIO
% ==============================================================================
\section{Cuestionario}
\label{sec:cuestionario}

\begin{enumerate}[leftmargin=1.5em,itemsep=0.04em,parsep=0pt,topsep=0.1em]

\item \textbf{Imposibilidad de transición $0 \to 1$ en evaluación:} En evaluación ($CLK=1$), el PMOS de precarga está apagado ($V_{GS,p}=0\,\text{V}$). El nodo dinámico $dyn$ es puramente capacitivo y carece de trayectoria hacia $V_{DD}$; por ende, no puede recargarse. Solo puede conservar $V_{DD}$ (PDN en corte) o descargarse a tierra ($1 \to 0$, PDN activa). El \textit{clock feedthrough} acopla sobreimpulsos transitorios, pero la celda es incapaz de generar una transición lógica ascendente $0 \to 1$.

\item \textbf{Capacitancia de entrada (dinámica vs estática):} En CMOS estático, la entrada conmuta compuertas NMOS y PMOS en paralelo ($C_{\text{in}} = C_{gn} + C_{gp}$); en dominó, la señal excita únicamente transistores NMOS ($C_{\text{in}} = C_{gn}$), asumiendo el reloj la precarga PMOS. En la NAND3 caracterizada ($W_n = 270\,\text{nm}, W_p = 135\,\text{nm}$), suprimir la rama PMOS reduce la capacitancia de entrada en SPICE un $42.4\%$ ($0.3417\,\text{fF}$ vs $0.5936\,\text{fF}$), aliviando significativamente el esfuerzo y retardo del excitador.

\item \textbf{Factor de actividad $\alpha \approx 1$ e implicancias en potencia:} El nodo $dyn$ se precarga a $V_{DD}$ cada ciclo ($CLK=0$). Si las entradas activan la PDN, $dyn$ se descarga en cada evaluación ($CLK=1$), imponiendo $\alpha = 1$ aun con datos constantes. Con baja actividad de datos (donde estático ahorra con $\alpha < 0.1$), la lógica dinámica disipa continuamente potencia de reloj y conmutación ($P = \alpha C_L (V_{DD})^2 f$), perdiendo competitividad energética.

\item \textbf{Deducción de $\Delta V$ y condición de corte por $V_{th}$:} Por balance de carga con nodo interno en cero: $Q_0 = C_L V_{DD} = (C_L+C_a)V_f \implies \Delta V = [C_a/(C_a+C_L)]V_{DD}$. Esta relación deja de ser válida si el NMOS se estrangula antes de igualar potenciales ($V_{GS} < V_{th}$), condición que ocurre cuando $V_{th,\text{eff}} > [C_a/(C_a+C_L)]V_{DD}$. Al cortarse en $V_a = V_{DD} - V_{th,\text{eff}}$, cesa la redistribución y la caída efectiva queda acotada por $\Delta V = (C_a/C_L)(V_{DD} - V_{th,\text{eff}})$.

\item \textbf{Orden de apilamiento en PDN dinámica vs estática:} En lógica estática, el nivel final $V_{OL}$ solo exige conducción DC a tierra, siendo indiferente el orden de apilamiento. En lógica dinámica, $dyn$ flota en alta impedancia; situar las entradas activas en transistores superiores contiguos a $dyn$ conecta de inmediato capacitancias de difusión parásitas no precargadas, maximizando el \textit{charge sharing}. Situar el transistor en corte arriba aísla la carga interna y preserva $V_{dyn}$.

\item \textbf{Regla de monotonicidad dominó y rol del inversor:} En evaluación ($CLK=1$), las entradas a la PDN solo deben realizar transiciones $0 \to 1$ (o reposar en 0). Como $dyn$ parte de $V_{DD}$ y solo puede caer ($1 \to 0$), conectarlo directo causaría la descarga espuria inmediata de la etapa siguiente antes de que evalúe. El inversor estático complementa la señal ($out: 0 \to 1$), manteniendo la PDN receptora apagada durante precarga y activándola monótonamente solo tras evaluar.

\item \textbf{Funciones no realizables en dominó y soluciones:} La lógica dominó estándar es no inversora ($dyn \downarrow \implies out \uparrow$), impidiendo realizar funciones complementarias directas (NAND, NOR, XOR) en una sola etapa. Se soluciona mediante: (1) \textit{Dual-Rail Domino} (DCVSL), que genera a la vez $F$ y $\overline{F}$ con PDNs diferenciales complementarias; y (2) \textit{NP-Domino} (Zipper), alternando etapas n-dominó y p-dominó con relojes en oposición.

\item \textbf{Corriente subumbral y mayor $nfactor$ en 45nm HP vs LP:} La corriente subumbral escala como $I_{\text{sub}} \propto (n-1)\exp[(V_{GS}-V_{th})/(n k T/q)][1 - \exp(-V_{DS}/(kT/q))]$. El factor $n = 1 + C_{\text{dep}}/C_{\text{ox}}$ refleja el acoplamiento capacitivo compuerta-canal. En 45nm HP, el dopaje \textit{halo} severo para mitigar DIBL eleva $C_{\text{dep}}$ frente a $C_{\text{ox}}$, incrementando $n$ ($1.53$ vs $1.15$ en LP). Un mayor $n$ degrada la pendiente subumbral ($S = n(kT/q)\ln 10$), disparando la fuga $I_{\text{off}}$.

\item \textbf{Comparativa 45nm HP vs LP en lógica dinámica:} LP presenta óxido más grueso ($t_{oxp}=1.50\,\text{nm}$; $t_{oxe}=1.80\,\text{nm}/1.82\,\text{nm}$, $d_{tox}=0.30/0.32\,\text{nm}$ en NMOS/PMOS vs $t_{oxp}=1.00\,\text{nm}$, $t_{oxe}=1.25/1.30\,\text{nm}$, $d_{tox}=0.25/0.30\,\text{nm}$ en HP), mayor umbral ($0.623\,\text{V}$ vs $0.469\,\text{V}$) y fuga $186\times$ menor ($0.67\,\text{pA}$ vs $125.1\,\text{pA}$), reteniendo carga a 85°C ($1\,\mu\text{s}$, $\Delta V < 2\,\text{mV}$) sin keeper. Empero, conmuta $52.5\times$ más lento: LP prima en retención y HP en velocidad.

\item \textbf{Keeper condicional y contención:} El keeper PMOS mitiga fugas y \textit{charge sharing}, pero disputa corriente con la PDN en evaluación, demorando la conmutación ($+12.90\%$ en 45nm HP con $W_{kp}=45\,\text{nm}$). El keeper condicional reduce esta penalización: retiene con un transistor mínimo en reposo y desconecta la rama mayor mediante reloj retardado durante la evaluación activa, reduciendo la demora a solo $+5.11\%$ y asegurando retención a 85°C.

\pagebreak
\item \textbf{Realimentación de keeper desde la salida del inversor:} El keeper PMOS restaura $V_{DD}$ en $dyn$; para encenderlo, su compuerta requiere $0\,\text{V}$. Conectar la compuerta directamente a $dyn$ lo mantendría apagado ($V_G=V_S=V_{DD}$). La salida invertida $out = \overline{dyn}$ provee la polaridad complementaria requerida. Asimismo, la alta ganancia regenerativa del inversor proporciona conmutación abrupta y desacopla la compuerta del keeper de perturbaciones transitorias en el nodo flotante.

\item \textbf{Clock feedthrough, origen en PMOS y dependencia con $dV/dt$:} El \textit{clock feedthrough} es el acoplamiento capacitivo del flanco de reloj al nodo flotante $dyn$ vía la capacitancia de solapamiento $C_{gd} = (c_{gdo}+c_{gdl})W_p$ del PMOS de precarga. La corriente inyectada es $i = C_{gd}(dV_{\text{clk}}/dt)$; flancos más abruptos (menor $t_r, t_f$) elevan $dV/dt$, inyectando un sobreimpulso ($+64.3\,\text{mV}$ en 45nm HP) que es recortado por conducción del diodo drenaje-sustrato $D_{DB}$ del PMOS.

\item \textbf{Skew positivo tolerable en footed vs cortocircuito en unfooted:} En dominó \textit{Footed}, el transistor de pie $M_f$ aísla la PDN hasta la llegada de su reloj local $CLK2$; si $T_{\text{skew}}>0$, la conmutación se retrasa pero sin degradar la lógica. En \textit{Unfooted}, la PDN conecta directo a masa; si la etapa previa inicia precarga ($CLK1: 1 \to 0$) mientras la receptora continúa en evaluación ($CLK2=1$), PMOS y PDN conducen simultáneamente, creando un cortocircuito $V_{DD} \to \text{GND}$ ($+94.1\%$ en corriente).

\item \textbf{Restricción de $f_{\min}$ en lógica dinámica vs estática:} El nodo dinámico almacena el estado como carga en un condensador flotante; las fugas subumbral y GIDL drenan esta carga en un tiempo finito $t_{\text{hold}}$, fijando una cota inferior de reloj $f \ge f_{\min} \approx 1/(2t_{\text{hold}})$ para evitar cruzar el umbral $V_M$. La lógica CMOS estática carece de límite inferior ($f_{\min}=0\,\text{Hz}$) porque sus salidas están conectadas permanentemente a los rieles de alimentación mediante trayectorias óhmicas.

\item \textbf{Degradación exponencial de $I_{\text{on}}/I_{\text{off}}$ en NTV:} Al reducir $V_{DD}$ hacia $V_{th}$, la corriente activa decae exponencialmente con la sobremarcha ($I_{\text{on}} \propto \exp[(V_{DD}-V_{th})/(n k T/q)]$), mientras la fuga $I_{\text{off}}$ se reduce marginalmente, colapsando el cociente $I_{\text{on}}/I_{\text{off}}$. La lógica dinámica es hipersensible a esto porque requiere $I_{\text{on}}$ alto para evaluar y $I_{\text{off}}$ ínfimo para retener, contrayendo el margen entre $t_{\text{eval}}$ y $t_{\text{hold}}$.

\item \textbf{Inviabilidad sin keeper a $0.4\,\text{V}$ y 85°C:} Aunque la ventana periódica abre ($t_{\text{hold}}=448\,\text{ns} \gg t_{\text{eval}}$), falla por dos causas: (1) Fugas térmicas degradan $dyn$ impidiendo pausas $\Delta t > 448\,\text{ns}$ ($V_{dyn} < V_{DD}/2$); y (2) Sin restauración activa colapsa $NM_H$, dejando al nodo vulnerable a ruido capacitivo.

\item \textbf{Imprescindibilidad de $AS/AD/PS/PD$ en BSIM4:} Los parámetros $AS, AD, PS, PD$ especifican áreas y perímetros de difusión para calcular las capacitancias de juntura parásitas ($C_j, C_{jsw}, C_{jswg}$). Al omitirse en SPICE, el simulador asigna cero ($C_{\text{diff}}=0$), modelando únicamente solapamientos de compuerta ($C_{gd}, C_{gs}$). Esto subestima la capacitancia parásita interna en $80\%\text{--}95\%$, enmascarando casi por completo el \textit{charge sharing} ($\Delta V \approx 0$ en vez de $164\,\text{mV}$).

\item \textbf{Necesidad de \texttt{gmin} y \texttt{abstol} en nodos flotantes:} En estática, las bajas impedancias ($k\Omega$) toleran valores estándar. En dominó, $dyn$ flota con impedancias de teraohmios ($T\Omega$). El parámetro por defecto $\texttt{gmin}=1\text{e-}12$ inserta una conductancia espuria que drena picoamperios artificiales, falseando la retención; calibrar a $\texttt{gmin}=1\text{e-}15$ y $\texttt{abstol}=1\text{e-}14$ asegura que las fugas medidas reflejen la física BSIM4 y no tolerancias del solver.

\item \textbf{Relación dominó con TGFF imperfecto (Bappy et al., SOCC 2026):} Las compuertas dominó exigen a las etapas secuenciales excitadoras niveles estables o transiciones monótonas ($0 \to 1$) en evaluación. El TGFF propuesto por Bappy, Zhao y Han (SOCC 2026 [5]) sustituye compuertas de transmisión balanceadas por transistores de paso para NTV; esto introduce caídas $V_{th}$ y rebotes capacitivos que pueden encender falsamente la PDN receptora o degradar la evaluación dominó.

\item \textbf{Efecto cuerpo en PDN apilada y viabilidad en bulk:} En transistores apilados, los nodos internos suben ($V_S > 0\,\text{V}$), polarizando el sustrato en reversa ($V_{SB}>0\,\text{V}$) y elevando el umbral: $\Delta V_{th} = \gamma(\sqrt{2\phi_F + V_{SB}} - \sqrt{2\phi_F})$, lo que merma la sobremarcha y retarda la evaluación. Conectar sustrato a fuente anularía el efecto, pero es inviable en procesos \textit{bulk} con pozo/sustrato común; mitigarlo exige triple-well o SOI.

\end{enumerate}

